% section: 二項定理

  \section{二項定理}
    次に,和の累乗を展開するための公式を確認する.
    自然数$n$と実数$x,y$について,
    \[
      (x+y)^n=\sum_{k=0}^{n}\binom{n}{k}x^{n-k}y^{k}
    \]
    が成り立つ.これを二項定理という.ここで$\dbinom{n}{k}$は二項係数であり,
    \[
      \binom{n}{k}=\frac{n!}{k!(n-k)!}
    \]
    で定義される. ここで,$k$は$0\le k\le n$を満たす整数であり,$0!=1$とする.
    この二項係数は,$n$個から$k$個を選ぶ組合せの数に一致する.
    例えば$n=2$では,
    \[
      (x+y)^2=x^2+2xy+y^2
    \]
    $n=3$では,
    \[
      (x+y)^3=x^3+3x^2y+3xy^2+y^3
    \]
    となる.
    後の節で扱う複素数に対しても,同じ二項定理が成り立つ.
    付録では,この公式を用いて差分を取ると多項式の次数が下がる理由などを説明する.